% GENERATED FILE. Edit canonical lessons, then run npm run generate:books.
% canonical-source-hash: fnv1a64:3a0e056757573ccc
% canonical-lessons: KA-C336-sainya, KA-C336-senapati, KA-C336-yodha, KA-C336-horadu, KA-C336-himmettu

\chapter{The Army Marches}
\label{ch:ka-the-army-marches}
\hlchaptermodality{336}

\begin{chapteropening}
\textbf{By the end of this chapter:} \emph{I can say an army, a commander, a warrior, to fight and to retreat.}

Complete the last lesson of chapter 336: I can say an army, a commander, a warrior, to fight and to retreat.
\end{chapteropening}

\section[\texorpdfstring{\kn{ಸೈನ್ಯ} (sainya)}{ಸೈನ್ಯ (sainya)}]{\kn{ಸೈನ್ಯ} (sainya) — an army}

\label{lesson:KA-C336-sainya}



\begin{quote}
Before the new one: say the Kannada for a leaflet, then the Kannada for a publisher.
\end{quote}

\subsection*{You'll want to know: \kn{ಸೈನ್ಯ}}
\textbf{\kn{ಸೈನ್ಯ}} — \emph{sainya} — \textquotedblleft{}an army\textquotedblright{}.

\begin{quote}
\kn{ರಾಜನ} \kn{ಸೈನ್ಯ} \kn{ನದಿಯನ್ನು} \kn{ದಾಟಿತು}.
\end{quote}

\begin{quote}
\emph{rājana sainya nadiyannu dāṭitu.}
\end{quote}

\begin{quote}
The king's army crossed the river.
\end{quote}

A soldier is a \kn{ಸೈನಿಕ}; the army is the \kn{ಸೈನ್ಯ}.

\begin{cousinweb}
\textbf{\kn{ಸೈನ್ಯ}} (\emph{sainya}) is borrowed from Sanskrit \emph{sainya}, \textquotedblleft{}belonging to an army\textquotedblright{}, and as a noun \textquotedblleft{}an army\textquotedblright{} (CDIAL 13602), from \emph{senā}, \textquotedblleft{}an army\textquotedblright{}.
\end{cousinweb}

\begin{grammarlens}[title={What you've built}]
One more word for telling what happened.
\end{grammarlens}

\subsection*{Your turn}
\begin{itemize}
\raggedright
  \item \emph{Say it:} \emph{sainya}
  \item \emph{Say it:} \emph{sainya}, once more
  \item \emph{Say it:} \emph{rājana sainya nadiyannu dāṭitu.} — The king's army crossed the river.
  \item \emph{Recall:} say the Kannada for a leaflet, then the Kannada for a publisher, then say \emph{sainya} again
\end{itemize}

\begin{tcolorbox}[breakable,colback=teal!4,colframe=teal!35!black,title={Before you move on}]
What does \kn{ಸೈನ್ಯ} mean? (\textquotedblleft{}An army\textquotedblright{}.) Say it once more.
\end{tcolorbox}
\section[\texorpdfstring{\kn{ಸೇನಾಪತಿ} (sēnāpati)}{ಸೇನಾಪತಿ (sēnāpati)}]{\kn{ಸೇನಾಪತಿ} (sēnāpati) — a commander; a general}

\label{lesson:KA-C336-senapati}



\begin{quote}
Before the new one: say the Kannada for a publisher, then the Kannada for an army.
\end{quote}

\subsection*{You'll want to know: \kn{ಸೇನಾಪತಿ}}
\textbf{\kn{ಸೇನಾಪತಿ}} — \emph{sēnāpati} — \textquotedblleft{}a commander; a general\textquotedblright{}.

\begin{quote}
\kn{ಸೇನಾಪತಿ} \kn{ಸೈನಿಕರಿಗೆ} \kn{ಮುಂದೆ} \kn{ಹೋಗಲು} \kn{ಹೇಳಿದನು}.
\end{quote}

\begin{quote}
\emph{sēnāpati sainikarige munde hōgalu hēḷidanu.}
\end{quote}

\begin{quote}
The general told the soldiers to go forward.
\end{quote}

Plural: \kn{ಸೇನಾಪತಿಗಳು}, \emph{sēnāpatigaḷu}.

\kn{ಸೇನೆ}, an army, and \kn{ಪತಿ}, a lord.

\begin{cousinweb}
\textbf{\kn{ಸೇನಾಪತಿ}} (\emph{sēnāpati}) is borrowed from Sanskrit \emph{senāpati}, \textquotedblleft{}the general of an army\textquotedblright{} (Monier-Williams), from \emph{senā} \textquotedblleft{}an army\textquotedblright{} and \emph{pati} \textquotedblleft{}a lord, a master\textquotedblright{}.
\end{cousinweb}

\begin{grammarlens}[title={What you've built}]
One more word for telling what happened.
\end{grammarlens}

\subsection*{Your turn}
\begin{itemize}
\raggedright
  \item \emph{Say it:} \emph{sēnāpati}
  \item \emph{Say it:} \emph{sēnāpati}, once more
  \item \emph{Say it:} \emph{sēnāpati sainikarige munde hōgalu hēḷidanu.} — The general told the soldiers to go forward.
  \item \emph{Recall:} say the Kannada for a publisher, then the Kannada for an army, then say \emph{sēnāpati} again
\end{itemize}

\begin{tcolorbox}[breakable,colback=teal!4,colframe=teal!35!black,title={Before you move on}]
What does \kn{ಸೇನಾಪತಿ} mean? (\textquotedblleft{}A commander; a general\textquotedblright{}.) Say it once more.
\end{tcolorbox}
\section[\texorpdfstring{\kn{ಯೋಧ} (yōdha)}{ಯೋಧ (yōdha)}]{\kn{ಯೋಧ} (yōdha) — a warrior}

\label{lesson:KA-C336-yodha}



\begin{quote}
Before the new one: say the Kannada for an army, then the Kannada for a commander.
\end{quote}

\subsection*{You'll want to know: \kn{ಯೋಧ}}
\textbf{\kn{ಯೋಧ}} — \emph{yōdha} — \textquotedblleft{}a warrior\textquotedblright{}.

\begin{quote}
\kn{ಆ} \kn{ಯೋಧ} \kn{ಯುದ್ಧದಲ್ಲಿ} \kn{ಊರನ್ನು} \kn{ಕಾಪಾಡಿದನು}.
\end{quote}

\begin{quote}
\emph{ā yōdha yuddhadalli ūrannu kāpāḍidanu.}
\end{quote}

\begin{quote}
That warrior protected the town in the war.
\end{quote}

Plural: \kn{ಯೋಧರು}, \emph{yōdharu}.

Today it is also a respectful word for a soldier: \kn{ನಮ್ಮ} \kn{ಯೋಧರು}.

\begin{cousinweb}
\textbf{\kn{ಯೋಧ}} (\emph{yōdha}) is borrowed from Sanskrit \emph{yodha}, \textquotedblleft{}a fighter, a warrior, a soldier\textquotedblright{} (Monier-Williams), from the root \emph{yudh}, \textquotedblleft{}to fight\textquotedblright{}, as in \kn{ಯುದ್ಧ}.
\end{cousinweb}

\begin{grammarlens}[title={What you've built}]
One more word for telling what happened.
\end{grammarlens}

\subsection*{Your turn}
\begin{itemize}
\raggedright
  \item \emph{Say it:} \emph{yōdha}
  \item \emph{Say it:} \emph{yōdha}, once more
  \item \emph{Say it:} \emph{ā yōdha yuddhadalli ūrannu kāpāḍidanu.} — That warrior protected the town in the war.
  \item \emph{Recall:} say the Kannada for an army, then the Kannada for a commander, then say \emph{yōdha} again
\end{itemize}

\begin{tcolorbox}[breakable,colback=teal!4,colframe=teal!35!black,title={Before you move on}]
What does \kn{ಯೋಧ} mean? (\textquotedblleft{}A warrior\textquotedblright{}.) Say it once more.
\end{tcolorbox}
\section[\texorpdfstring{\kn{ಹೋರಾಡು} (hōrāḍu)}{ಹೋರಾಡು (hōrāḍu)}]{\kn{ಹೋರಾಡು} (hōrāḍu) — to fight; to struggle}

\label{lesson:KA-C336-horadu}



\begin{quote}
Before the new one: say the Kannada for a commander, then the Kannada for a warrior.
\end{quote}

\subsection*{You'll want to know: \kn{ಹೋರಾಡು}}
\textbf{\kn{ಹೋರಾಡು}} — \emph{hōrāḍu} — \textquotedblleft{}to fight; to struggle\textquotedblright{}.

\begin{quote}
\kn{ರೈತರು} \kn{ತಮ್ಮ} \kn{ನೆಲಕ್ಕಾಗಿ} \kn{ಹೋರಾಡಿದರು}.
\end{quote}

\begin{quote}
\emph{raitaru tamma nelakkāgi hōrāḍidaru.}
\end{quote}

\begin{quote}
The farmers fought for their land.
\end{quote}

Past (I …): \kn{ಹೋರಾಡಿದೆ}, \emph{hōrāḍide}.

What one fights for takes -\kn{ಆಗಿ}: \kn{ನ್ಯಾಯಕ್ಕಾಗಿ} \kn{ಹೋರಾಡು}, to fight for justice.

\begin{cousinweb}
\textbf{\kn{ಹೋರಾಡು}} (\emph{hōrāḍu}) is a native Dravidian compound: \emph{pōr}, \textquotedblleft{}to fight, to wrestle\textquotedblright{} (DEDR 4540), whose \emph{p-} became \emph{h-}, and \emph{āḍu}, \textquotedblleft{}to move, to act\textquotedblright{} (DEDR 347).
\end{cousinweb}

\begin{grammarlens}[title={What you've built}]
One more word for telling what happened.
\end{grammarlens}

\subsection*{Your turn}
\begin{itemize}
\raggedright
  \item \emph{Say it:} \emph{hōrāḍu}
  \item \emph{Say it:} \emph{hōrāḍu}, once more
  \item \emph{Say it:} \emph{raitaru tamma nelakkāgi hōrāḍidaru.} — The farmers fought for their land.
  \item \emph{Recall:} say the Kannada for a commander, then the Kannada for a warrior, then say \emph{hōrāḍu} again
\end{itemize}

\begin{tcolorbox}[breakable,colback=teal!4,colframe=teal!35!black,title={Before you move on}]
What does \kn{ಹೋರಾಡು} mean? (\textquotedblleft{}To fight; to struggle\textquotedblright{}.) Say it once more.
\end{tcolorbox}
\section[\texorpdfstring{\kn{ಹಿಮ್ಮೆಟ್ಟು} (himmeṭṭu)}{ಹಿಮ್ಮೆಟ್ಟು (himmeṭṭu)}]{\kn{ಹಿಮ್ಮೆಟ್ಟು} (himmeṭṭu) — to retreat; to fall back}

\label{lesson:KA-C336-himmettu}



\begin{quote}
Before the new one: say the Kannada for a warrior, then the Kannada for to fight.
\end{quote}

\subsection*{You'll want to know: \kn{ಹಿಮ್ಮೆಟ್ಟು}}
\textbf{\kn{ಹಿಮ್ಮೆಟ್ಟು}} — \emph{himmeṭṭu} — \textquotedblleft{}to retreat; to fall back\textquotedblright{}.

\begin{quote}
\kn{ದೊಡ್ಡ} \kn{ಸೈನ್ಯ} \kn{ಬರುವುದನ್ನು} \kn{ನೋಡಿ} \kn{ಅವರು} \kn{ಹಿಮ್ಮೆಟ್ಟಿದರು}.
\end{quote}

\begin{quote}
\emph{doḍḍa sainya baruvudannu nōḍi avaru himmeṭṭidaru.}
\end{quote}

\begin{quote}
Seeing a large army coming, they fell back.
\end{quote}

Past (I …): \kn{ಹಿಮ್ಮೆಟ್ಟಿದೆ}, \emph{himmeṭṭide}.

\kn{ಹಿಂ}, back, and \kn{ಮೆಟ್ಟು}, to step: to step back.

\begin{cousinweb}
\textbf{\kn{ಹಿಮ್ಮೆಟ್ಟು}} (\emph{himmeṭṭu}) is a native Dravidian compound: \emph{him}, \textquotedblleft{}behind, back\textquotedblright{} (DEDR 4205), and \emph{meṭṭu}, \textquotedblleft{}to put the foot down, to step\textquotedblright{} (DEDR 5057).
\end{cousinweb}

\begin{grammarlens}[title={What you've built}]
One more word for telling what happened.
\end{grammarlens}

\subsection*{Your turn}
\begin{itemize}
\raggedright
  \item \emph{Say it:} \emph{himmeṭṭu}
  \item \emph{Say it:} \emph{himmeṭṭu}, once more
  \item \emph{Say it:} \emph{doḍḍa sainya baruvudannu nōḍi avaru himmeṭṭidaru.} — Seeing a large army coming, they fell back.
  \item \emph{Recall:} say the Kannada for a warrior, then the Kannada for to fight, then say \emph{himmeṭṭu} again
\end{itemize}

\begin{tcolorbox}[breakable,colback=teal!4,colframe=teal!35!black,title={Before you move on}]
What does \kn{ಹಿಮ್ಮೆಟ್ಟು} mean? (\textquotedblleft{}To retreat; to fall back\textquotedblright{}.) Say it once more.
\end{tcolorbox}
